% Figure: ME/CFS Modeling Framework Overview
% Shows modeling approaches, subsystem models, and predictive applications

\begin{figure}[htbp]
\centering
\begin{tikzpicture}[
    scale=0.78, every node/.style={scale=0.78},
    % Approach boxes (top row, blue tones)
    approach/.style={draw=blue!70!black, fill=blue!20, very thick, rounded corners,
        text width=3cm, align=center, minimum height=1.2cm, font=\small\bfseries},
    % Subsystem boxes (middle row, teal/green tones)
    subsystem/.style={draw=teal!70!black, fill=teal!15, very thick, rounded corners,
        text width=3.2cm, align=center, minimum height=1.4cm, font=\small\bfseries},
    % Application boxes (bottom row, orange tones)
    application/.style={draw=orange!70!black, fill=orange!15, very thick, rounded corners,
        text width=4.5cm, align=center, minimum height=1.2cm, font=\small\bfseries},
    % Arrow styles
    feeds/.style={-latex, very thick, line width=1.2pt, blue!60!black},
    couples/.style={-latex, thick, line width=1pt, teal!60!black, dashed},
    integrates/.style={-latex, very thick, line width=1.2pt, orange!70!black},
    % Annotation style
    annot/.style={font=\scriptsize\itshape, blue!60!black},
    coupling-label/.style={font=\scriptsize, teal!70!black, fill=white, inner sep=2pt,
        rounded corners=1pt},
]

% === TITLE ===
\node[font=\large\bfseries, blue!70!black] (title) at (7.5, 0)
    {ME/CFS Modeling Framework};

% === TOP ROW: Modeling Approaches ===
% 5cm spacing between centers for breathing room
\node[approach] (ode) at (0, -1.8)
    {Ordinary\\Differential\\Equations};
\node[approach] (stoch) at (5, -1.8)
    {Stochastic\\Models};
\node[approach] (net) at (10, -1.8)
    {Network\\Models};
\node[approach] (abm) at (15, -1.8)
    {Agent-Based\\Models};

% Section annotations for approaches
\node[annot, below=2pt of ode] {\S\,26.2};
\node[annot, below=2pt of stoch] {\S\,26.3};
\node[annot, below=2pt of net] {\S\,26.4};
\node[annot, below=2pt of abm] {\S\,26.5};

% === MIDDLE ROW: Subsystem Models ===
% Same 5cm spacing, aligned with approaches
\node[subsystem] (energy) at (0, -6.2)
    {Energy\\Metabolism\\{\footnotesize Ch.\,27}};
\node[subsystem] (immune) at (5, -6.2)
    {Immune\\System\\{\footnotesize Ch.\,28}};
\node[subsystem] (neuro) at (10, -6.2)
    {Neuroendocrine\\System\\{\footnotesize Ch.\,29}};
\node[subsystem] (integrated) at (15, -6.2)
    {Integrated\\Systems\\{\footnotesize Ch.\,30}};

% === BOTTOM ROW: Applications ===
\node[application] (temporal) at (4, -10.5)
    {Temporal Evolution\\{\footnotesize Ch.\,31}};
\node[application] (predict) at (11.2, -10.5)
    {Predictive Applications\\{\footnotesize Ch.\,32}};

% === ARROWS: Approaches → Subsystems ===
% Each approach drops to its aligned subsystem (direct),
% plus one lateral routed connection

% ODE → Energy (direct), ODE → Immune (routed)
\draw[feeds] (ode.south) -- (energy.north);
\draw[feeds] (ode.south) -- ++(0,-0.8) -| ([xshift=-4pt]immune.north);

% Stochastic → Immune (direct), Stochastic → Neuro (routed)
\draw[feeds] (stoch.south) -- (immune.north);
\draw[feeds] (stoch.south) -- ++(0,-1.0) -| ([xshift=-4pt]neuro.north);

% Network → Neuro (direct), Network → Integrated (routed)
\draw[feeds] (net.south) -- (neuro.north);
\draw[feeds] (net.south) -- ++(0,-0.8) -| ([xshift=-4pt]integrated.north);

% ABM → Integrated (direct)
\draw[feeds] (abm.south) -- (integrated.north);

% === DASHED ARROWS: Cross-system couplings ===
% With 5cm spacing there is ~1.5cm gap between node edges for labels
\draw[couples, bend left=18] (energy.east) to
    node[coupling-label, above=-1pt, pos=0.5] {cytokine--metabolic} (immune.west);
\draw[couples, bend left=18] (immune.east) to
    node[coupling-label, above=-1pt, pos=0.5] {neuroimmune} (neuro.west);
\draw[couples, bend left=18] (neuro.east) to
    node[coupling-label, above=-1pt, pos=0.5] {multi-system} (integrated.west);

% Long-range coupling: Energy ↔ Neuro (below the row)
\draw[couples, bend right=18] ([yshift=-2pt]energy.south east) to
    node[coupling-label, below, pos=0.5] {HPA--metabolic axis}
    ([yshift=-2pt]neuro.south west);

% === ARROWS: Subsystems → Applications ===
% Energy and Immune feed Temporal
\draw[integrates] (energy.south) -- ++(0,-0.5) -| ([xshift=-8pt]temporal.north);
\draw[integrates] (immune.south) -- ++(0,-0.7) -| ([xshift=8pt]temporal.north);

% Neuro and Integrated feed Predictive
\draw[integrates] (neuro.south) -- ++(0,-0.5) -| ([xshift=-8pt]predict.north);
\draw[integrates] (integrated.south) -- ++(0,-0.7) -| ([xshift=8pt]predict.north);

% Cross-connections (lighter): Integrated → Temporal, Energy → Predictive
\draw[integrates, opacity=0.4] (integrated.south) -- ++(0,-0.9) -|
    ([xshift=2pt]temporal.north east);
\draw[integrates, opacity=0.4] (energy.south) -- ++(0,-0.9) -|
    ([xshift=-2pt]predict.north west);

% Temporal → Predictive
\draw[integrates] (temporal.east) --
    node[below, font=\scriptsize, orange!70!black] {trajectory forecasting}
    (predict.west);

% === LEGEND ===
\node[font=\footnotesize, anchor=north west] at (0.5, -12.3) {%
    \begin{tikzpicture}[baseline=-0.5ex]
        \draw[feeds] (0,0) -- (0.8,0);
    \end{tikzpicture}~Methodological basis \qquad
    \begin{tikzpicture}[baseline=-0.5ex]
        \draw[couples] (0,0) -- (0.8,0);
    \end{tikzpicture}~Cross-system coupling \qquad
    \begin{tikzpicture}[baseline=-0.5ex]
        \draw[integrates] (0,0) -- (0.8,0);
    \end{tikzpicture}~Integration%
};

\end{tikzpicture}
\caption{Overview of the ME/CFS modeling framework.
Modeling approaches (top) are applied to physiological subsystems (middle),
which couple into integrated and predictive models (bottom).
Dashed arrows indicate cross-system couplings.}
\label{fig:modeling-framework-overview}
\end{figure}
